%\documentclass[12pt,oneside]{amsart}
\documentclass{scrartcl}

%\documentclass{article}

\usepackage{amsthm,amsmath,amssymb}
\usepackage[numbers]{natbib}

%\usepackage[colorlinks]{hyperref}
\usepackage{hyperref}
\hypersetup{
    colorlinks,%
    citecolor=black,%
    filecolor=black,%
    linkcolor=black,%
    urlcolor=black
}
\usepackage{hypernat}
\usepackage[top=2.5cm, bottom=2.5cm, left=2.8cm,
right=2.8cm]{geometry}
\usepackage{graphicx}


%\setlength{\parskip}{0pt}
%\setlength{\parsep}{0pt}
%\setlength{\headsep}{0pt}
%\setlength{\topskip}{0pt}
%\setlength{\topmargin}{0pt}
%\setlength{\topsep}{0pt}
%\setlength{\partopsep}{0pt}

\linespread{1}

%\usepackage{marvosym}

%\textwidth = 6.5 in \textheight = 9.2 in \oddsidemargin = 0.0 in
%\evensidemargin = 0.0 in \topmargin = -0.0 in \headheight = 0.0 in
%\headsep = 0.0in \parskip = 0.0in \parindent = 0.0in
%\def\baselinestretch{0.871}


\newtheorem{theorem}{Theorem}[section]
\newtheorem{proposition}[theorem]{Proposition}
\newtheorem{problem}[theorem]{Problem}
\newtheorem{fact}[theorem]{Fact}
\newtheorem{lemma}[theorem]{Lemma}
\newtheorem{defn}[theorem]{Definition}
\newtheorem{corollary}[theorem]{Corollary}
\newtheorem{remark}{Remark}[section]
\newtheorem{example}[theorem]{Example}

\newcommand{\E}{{\mathbb E}}
\newcommand{\se}{{\mathcal{E}}}
\newcommand{\R}{{\mathbb R}}
\renewcommand{\P}{{\mathbb P}}
\newcommand{\ac}{{\mathcal{A}}}
\newcommand{\A}{{\mathcal{A}}}
\newcommand{\B}{{\mathcal{B}}}
\newcommand{\C}{{\mathcal{C}}}
\newcommand{\M}{{\mathcal{M}}}
\newcommand{\Mf}{{\mathfrak{M}}}
\newcommand{\T}{{\mathcal{T}}}
\newcommand{\Z}{{\mathcal{Z}}}
\newcommand{\K}{{\mathcal{K}}}
\newcommand{\lc}{{\mathcal{L}}}
\newcommand{\Ps}{{\mathcal{P}}}
\newcommand{\G}{{\mathcal{G}}}
\newcommand{\pa}{{\mathring{p}}}
\newcommand{\F}{{\cal F}}
\newcommand{\samp}{{\mathcal{X}}}
\newcommand{\X}{{\mathcal{X}}}
\newcommand{\hi}{{\hat{\Phi}}}
\newcommand{\data}{{\text{data}}}
\newcommand{\relu}{{\text{ReLU}}}

\newcommand{\ridge}{{\hat{\beta}_{\text{ridge}}(\lambda)}}
\newcommand{\lasso}{{\hat{\beta}_{\text{lasso}}(\lambda)}}
\newcommand{\ridgemu}{{\hat{\mu}_t^{\text{ridge}}(\lambda)}}
\newcommand{\lassomu}{{\hat{\mu}_t^{\text{lasso}}(\lambda)}}


\newcounter{rcnt}[section]
\renewcommand{\thercnt}{(\roman{rcnt})}

\def\qt#1{\qquad\text{#1}}

\def\argmin{\mathop{\rm argmin}}
\def\argmax{\mathop{\rm argmax}}


\setlength{\parskip}{1.4 \medskipamount}
%\sloppy
%\linespread{1.3}

\begin{document}

\title{STAT 238 - Bayesian Statistics  \\
  Lecture Thirty Five} 
\subtitle{Spring 2026, UC Berkeley}

\author{Aditya Guntuboyina}


\date{24 April 2026}

\maketitle

\section{HMC and Leapfrog Discretization}

There is a given target probability density $\pi$ that we need to
sample from. In HMC, given a current state $x$, one attempts to solve the
dynamics:
\begin{align}\label{hmc_form}
    \begin{pmatrix}
    \dot{x}(t) \\ \dot{v}(t)
  \end{pmatrix} =
  \begin{pmatrix}
    \frac{\partial H}{\partial v}(x(t), v(t)) \\ -\frac{\partial
    H}{\partial x}(x(t), v(t))
  \end{pmatrix} ~\text{with initialization}~
  \begin{pmatrix}
    x(0) \\ v(0)
  \end{pmatrix} =
  \begin{pmatrix}
    x \\ v
  \end{pmatrix}
\end{align}
with $v \sim N(0, I_d)$. Here $H(\cdot, \cdot)$ denotes the
Hamiltonian:
\begin{align}\label{hamiltonian}
  H(x, v) = H(x_1, \dots, x_d, v_1, \dots, v_d) := -\log \pi(x) +
  \frac{1}{2}\|v\|^2. 
\end{align}
Let $T_{\sigma}$ denote the function which maps $(x, v)$ to
$(x(\sigma), v(\sigma))$ obtained by solving \eqref{hmc_form} for $t
\in [0, \sigma]$. We saw in the last lecture that $T_{\sigma}$ is
invertible (specifically, $T_{\sigma} S = (T_{\sigma} S)^{-1}$ where
$S$ is the velocity sign-flip operator), Hamiltonian preserving (i.e.,
$H(T_{\sigma}(x, v)) = H(x, v)$) and Volume preserving (i.e., the
determinant of the Jacobian of $T_{\sigma}$ always equals 1).

The HMC ODE \eqref{hmc_form} cannot be solved exactly in practice, so
one attempts to solve it approximately via discretization by fixing a
step size $\epsilon$ and solving \eqref{hmc_form} approximately for $t
= 0, \epsilon, 2 \epsilon, \dots$. To go from $(x(t), v(t))$ to
$(x(t+\epsilon), v(t + \epsilon))$, one uses the leapfrog
discretization:
\begin{equation}
  \label{leapfrog}
  \begin{split}
&    v\left(t + \frac{\epsilon}{2}\right) = v(t) + \frac{\epsilon}{2}
  \nabla \log \pi(x(t)) \\
& x(t + \epsilon) = x(t) + \epsilon v\left(t + \frac{\epsilon}{2}
  \right) \\
& v(t + \epsilon) = v \left(t + \frac{\epsilon}{2} \right) +
  \frac{\epsilon}{2} \nabla \log \pi(x(t + \epsilon))
  \end{split}
\end{equation}
If we denote this map from $(x(t), y(t))$ to $(x(t + \epsilon), y(t +
\epsilon))$ by $T^{\text{disc}}_{\epsilon}(x, v)$ and then apply this
mapping $N$ times in succession starting from $(x, v)$, then the
overall mapping is given by
\begin{align*}
  T^{\text{disc}, N}_{\epsilon} = T_{\epsilon}^{\text{disc}} \circ \dots \circ
  T_{\epsilon}^{\text{disc}}. 
\end{align*}
The above mapping is an approximation to the continuous ODE flow
$T_{\sigma}$:
\begin{align*}
  T_{\epsilon}^{\text{disc}, N}(x, v) \approx T_{\sigma}(x, v) ~~
  \text{ provided $\epsilon = \frac{\sigma}{N}$}. 
\end{align*}
In the last lecture, we observed that $T_{\epsilon}^{\text{disc}, N}$
is also invertible (with $S T_{\epsilon}^{\text{disc}, N} =
(S T_{\epsilon}^{\text{disc}, N})^{-1}$) and volume preserving. But it
does not preserve the Hamiltonian i.e., $H(T_{\epsilon}^{\text{disc},
  N}(x, v))$ can be different from $H(x, v)$.

While implementing leapfrog discretization to compute $x(N\epsilon),
v(N \epsilon)$, it can be checked that the following iterations can be
used which combine the two velocity updates into one:
\begin{align*}
 & x(j \epsilon) = x((j-1)\epsilon) + \epsilon v \left((j - 1/2)
   \epsilon \right) \\
 & v((j+ 1/2) \epsilon) = v((j-1/2) \epsilon) + \epsilon x(j \epsilon)
\end{align*}
for $j = 1, \dots, N-1$. This iteration needs to be initialized using
$x(0) = x$ and $v(\epsilon/2) = v(0) + (\epsilon/2) \nabla \log
\pi(x(0))$. At iteration $N-1$, we obtain $x((N-1)\epsilon)$ and
$v((N-1/2) \epsilon)$. From here $x(N \epsilon)$ and $v(N \epsilon)$
can be computed by:
\begin{align*}
  &  x(N \epsilon) = x((N-1) \epsilon) + \epsilon v((N-1/2) \epsilon) \\
  & v(N \epsilon) = v((N-1/2)\epsilon) + (\epsilon/2) \nabla \log
    \pi(x(N \epsilon)). 
\end{align*}

\section{The HMC Algorithm}

We are now ready to formally describe the HMC algorithm. The algorithm
starts with an arbitrary initialization $x^{(0)}$ and then repeats the
following for $t = 0, 1, 2, \dots$:
\begin{enumerate}
\item Generate $v^{(t)} \sim N(0, I_d)$.
\item Run discretized Leapfrog Hamiltonian dynamics starting at
  $(x^{(t)}, v^{(t)})$ for some fixed step size $\epsilon$ and
  trajectory length $\sigma$. This gives $(y, w) =
  T_{\epsilon}^{\text{disc}, N}(x, v)$ (here $N = \sigma/\epsilon$ is
  the number of leapfrog steps).
\item Compute the acceptance probability: 
  \begin{align*}
    \min \left(1, \exp(H(x^{(t)}, v^{(t)}) - H(y, w)) \right)
  \end{align*}
  and update $x^{(t+1)}$ as:
  \[
x^{(t+1)} =
\begin{cases}
y & \text{with probability }     \min \left(1, \exp(H(x^{(t)},
     v^{(t)}) - H(y, w)) \right), \\[1em] 
x^{(t)} & \text{with probability } 1 - \min \left(1, \exp(H(x^{(t)},
     v^{(t)}) - H(y, w)) \right). 
\end{cases}
\]
\end{enumerate}
We need to verify that the Markov Chain given by the above algorithm
maintains detailed balance with respect to $\pi$. This is difficult to
verify directly for the chain with transitions to $x^{(t+1)}$ from
$x^{(t)}$. Instead, we shall verify detailed balance for the chain
which transitions to $(x^{(t+1)}, v^{(t+1)})$ from $(x^{(t)},
v^{(t)})$. Note that $v^{(t+1)}$ does not appear in the above
description of the HMC algorithm. We rewrite the algorithm below using
$v^{(t+1)}$ explicitly.

In addition to using $v^{(t+1)}$ explicitly, another change in the
following algorithm is that the equation $(y, w) =
T_{\epsilon}^{\text{disc}, N}(x, v)$ is now changed to $(y, w) = S
T_{\epsilon}^{\text{disc}, N}(x, v)$. In other words, we do a velocity
flip after running the discretized Hamiltonian dynamics. This velocity
flip makes the deterministic transition from $(x, v)$ to $(y, w)$ an
involution (i.e., its inverse is itself: $(S
T_{\epsilon}^{\text{disc}, N})^{-1} = S T_{\epsilon}^{\text{disc},
  N}$) which makes the argument for detailed balance very clean. Note
also that the velocity flip does not change the Hamiltonian (because
$H(y, w) = H(y, -w)$ as the dependence on $w$ is through $\|w\|^2$) so
the acceptance probability does not change. Algorithmically, both the
versions of the algorithm lead to identical output.

\begin{enumerate}
\item Generate $v^{(t)} \sim N(0, I_d)$.
\item Run discretized Leapfrog Hamiltonian dynamics (followed  by a
  velocity flip) starting at
  $(x^{(t)}, v^{(t)})$ for some fixed step size $\epsilon$ and
  trajectory length $\sigma$. This gives $(y, w) =
  S T_{\epsilon}^{\text{disc}, N}(x, v)$ (here $N = \sigma/\epsilon$ is
  the number of leapfrog steps).
\item Compute the acceptance probability: 
  \begin{align*}
    \min \left(1, \exp(H(x^{(t)}, v^{(t)}) - H(y, w)) \right)
  \end{align*}
  and update $(x^{(t+1)}, v^{(t+1)})$ as:
  \[
(x^{(t+1)}, v^{(t+1)}) =
\begin{cases}
(y, w) & \text{with probability }     \min \left(1, \exp(H(x^{(t)},
     v^{(t)}) - H(y, w)) \right), \\[1em] 
(x^{(t)} v^{(t)}), & \text{with probability } 1 - \min \left(1, \exp(H(x^{(t)},
     v^{(t)}) - H(y, w)) \right). 
\end{cases}
\] 
\end{enumerate}
Next we will verify detailed balance of the above Markov Chain with
respect to the joint density:
\begin{align*}
  \tilde{\pi}(x, v) \propto \exp \left(-H(x, v) \right).  
\end{align*}
Note that, under $\tilde{\pi}$, the variables $x \sim \pi$ and $v \sim
N(0, I_d)$ are independent (because $H(x, v) = -\log \pi(x) +
\|v\|^2/2$.

The transitions of the HMC Markov Chain from $(x, v)$ to $(y, w)$ are
deterministic (given by the involutive map $S
T_{\epsilon}^{\text{disc}, N}$) so we need to be careful about the
form of detailed balance.

\section{Recap: detailed balance}
So far in this course (e.g., see Lectures 25 and 31), we have used the
following definition of detailed balance. The Markov Chain with
transition densities given by $Q(x, y)$ satisfies detailed balance
with respect to the density $\pi(x)$ provided:
\begin{align}\label{detbal_cont}
  \pi(x) Q(x, y) = \pi(y) Q(y, x). 
\end{align}
This formulation is only true if the transitions have a density given
by $Q(x, \cdot)$. For HMC (viewed as a chain over $(x^{(t)},
v^{(t)})$), the transitions are deterministic so we cannot use the
detailed balance characterization in \eqref{detbal_cont}. The more
general way of writing the detailed balance condition is:
\begin{align}\label{gendetbal}
  \P\{\Theta^{(t)} \in A, \Theta^{(t+1)} \in B\} =   \P\{\Theta^{(t)}
  \in B, \Theta^{(t+1)} \in A\}
\end{align}
assuming that $\Theta^{(t)} \in A$ and the conditional distribution of
$\Theta^{(t+1)}$ given $\Theta^{(t)}$ is given by the Markov
chain. Also \eqref{gendetbal} is required to hold for all subsets $A$
and $B$ of the state space.

In the case when the Markov chain transitions are given by densities
$Q(x, \cdot)$, then it is easy to show that \eqref{gendetbal} is
equivalent to \eqref{detbal_cont}. For example, if we restrict $A$ to
be a small region around a point $x$, and $B$ to be a small region
around a point $y$, then the left hand side of \eqref{gendetbal} is
approximately $\pi(x) Q(x, y) \text{vol}(A) \text{vol}(B)$ and the
right hand side is $\pi(y) Q(y, x) \text{vol}(A) \text{vol}(B)$. One
can then cancel the product of the volume terms to obtain
\eqref{detbal_cont} from \eqref{gendetbal}.

On the other hand, when the Markov chain transitions do not have a
density, we need to work with the general condition
\eqref{gendetbal}. 

\section{Detailed Balance for HMC}
We shall now derive the detailed balance condition \eqref{gendetbal}
for HMC. We need to prove that:
\begin{align}\label{hmc_detbal}
  \P \left\{(x^{(t)}, v^{(t)}) \in A, (x^{(t+1)}, v^{(t+1)}) \in B
  \right\} =   \P \left\{(x^{(t)}, v^{(t)}) \in B, (x^{(t+1)},
  v^{(t+1)}) \in A 
  \right\}
\end{align}
where $(x^{(t)}, v^{(t)})$ has density proportional to $\exp(-H(x,
v))$. We compute the left hand side as
\begin{align*}
 & \P \left\{(x^{(t)}, v^{(t)}) \in A, (x^{(t+1)}, v^{(t+1)}) \in B
  \right\} \\ &= \int_A \exp(-H(x, v)) \P \left\{(x^{(t+1)}, v^{(t+1)})
             \in B \mid x^{(t)} = x, v^{(t)} = v \right\} dx dv
\end{align*}
To write the conditional probability $\P \left\{(x^{(t+1)}, v^{(t+1)})
             \in B \mid x^{(t)} = x, v^{(t)} = v \right\}$, note that
           $(x^{(t+1)}, v^{(t+1)})$ either equals $S
           T_{\epsilon}^{\text{disc}, N}(x, v)$ with probability
             $\min(1, \exp(H(x^{(t)}, v^{(t)}) - H(S
           T_{\epsilon}^{\text{disc}, N}(x, v)))$ or it equals $(x,
             v)$ with the opposite probability. For simplicity of
             notation, let us denote the function $S
             T_{\epsilon}^{\text{disc}, N}$ by $\Upsilon$. Thus
\begin{align*}
&  \P \left\{(x^{(t+1)}, v^{(t+1)})
  \in B \mid x^{(t)} = x, v^{(t)} = v \right\} \\
&= \min \left(1,\frac{\exp(-H(\Upsilon(x, v)))}{\exp(-H(x, v))}
  \right)  I\{\Upsilon(x, v) \in B\} + \left[1 - \min
  \left(1,\frac{\exp(-H(\Upsilon(x, v)))}{\exp(-H(x, v))} 
  \right) \right] I\{(x, v) \in B\}. 
\end{align*}

             We then get
\begin{align*}
 & \P \left\{(x^{(t)}, v^{(t)}) \in A, (x^{(t+1)}, v^{(t+1)}) \in B
  \right\} \\ &= \int_A \exp(-H(x, v)) \P \left\{(x^{(t+1)}, v^{(t+1)})
                \in B \mid x^{(t)} = x, v^{(t)} = v \right\} dx dv \\
  &= \int_A \exp(-H(x, v)) \min \left(1,\frac{\exp(-H(\Upsilon(x, v)))}{\exp(-H(x, v))}
    \right)  I\{\Upsilon(x, v) \in B\} dx dv \\
  &+ \int_A \exp(-H(x, v)) \left[1 - \min
  \left(1,\frac{\exp(-H(\Upsilon(x, v)))}{\exp(-H(x, v))} 
    \right) \right] I\{(x, v) \in B\} dx dv \\
  &= \int_A \min \left(\exp(-H(x, v)), \exp(-H(\Upsilon(x, v)))
    \right) I\{\Upsilon(x, v) \in B\} dx dv \\
  &+ \int_{A \cap B} \left[\exp(-H(x, v)) - \min \left(\exp(-H(x, v)), \exp(-H(\Upsilon(x, v)))
    \right)  \right] dx dv.
\end{align*}
By the same argument (but with $A$ and $B$ switched), the right hand
side of \eqref{hmc_detbal} satisfies: 
\begin{align*}
  &\P \left\{(x^{(t)}, v^{(t)}) \in B, (x^{(t+1)},
    v^{(t+1)}) \in A  \right\} \\
  &= \int_B \min \left(\exp(-H(x, v)), \exp(-H(\Upsilon(x, v)))
    \right) I\{\Upsilon(x, v) \in A\} dx dv \\
  &+ \int_{B \cap A} \left[\exp(-H(x, v)) - \min \left(\exp(-H(x, v)), \exp(-H(\Upsilon(x, v)))
    \right)  \right] dx dv.
\end{align*}
It is easy to see that the second terms in the expressions for $\P \left\{(x^{(t)}, v^{(t)}) \in A, (x^{(t+1)}, v^{(t+1)}) \in B
  \right\}$ and $\P \left\{(x^{(t)}, v^{(t)}) \in B, (x^{(t+1)},
    v^{(t+1)}) \in A  \right\}$ above are identical. So we only need to
  prove that the first terms also coincide i.e., we need to show
  \begin{equation}\label{hmc_tp}
    \begin{split}
   &\int_A \min \left(\exp(-H(x, v)), \exp(-H(\Upsilon(x, v)))
    \right) I\{\Upsilon(x, v) \in B\} dx dv \\ &=  \int_B \min \left(\exp(-H(x, v)), \exp(-H(\Upsilon(x, v)))
                                                 \right) I\{\Upsilon(x, v) \in A\} dx dv
   \end{split}                                              
\end{equation}
for all $A$ and $B$. For this, let us take the right hand side
integral and use the change of variable $(y, w) = \Upsilon(x, v)$, or
equivalently $(x, v) = \Upsilon^{-1}(y, w)$. Because $\Upsilon =
ST_{\epsilon}^{\text{disc}, N}$ is an involution, we have
$\Upsilon^{-1} = \Upsilon$. Thus
\begin{align*}
&  \int_B \min \left(\exp(-H(x, v)), \exp(-H(\Upsilon(x, v)))
  \right) I\{\Upsilon(x, v) \in A\} dx dv \\
&= \int \min \left(\exp(-H(x, v)), \exp(-H(\Upsilon(x, v)))
  \right) I\{\Upsilon(x, v) \in A, (x, v) \in B\} dx dv \\
&= \int \min \left(\exp(-H(y, w)), \exp(-H(\Upsilon(y, w)))
  \right) I\{\Upsilon(y, w) \in A, (y, w) \in B\} dy dw |\det J
  \Upsilon(y, w)|. 
\end{align*}
Because $\Upsilon$ is volume-preserving (we proved, in the last
lecture, that leapfrog discretization preserves volumes), the
determinant term above equals one. We thus have
\begin{align*}
&  \int_B \min \left(\exp(-H(x, v)), \exp(-H(\Upsilon(x, v)))
  \right) I\{\Upsilon(x, v) \in A\} dx dv \\
&= \int \min \left(\exp(-H(y, w)), \exp(-H(\Upsilon(y, w)))
  \right) I\{\Upsilon(y, w) \in A, (y, w) \in B\} dy dw  
\end{align*}
Similarly, one can also prove that the left hand side of
\eqref{hmc_tp} also equals the above. This completes the proof of
\eqref{hmc_tp} establishing detailed balance of the HMC markov chain.

For more details behind HMC, the paper \citet{neal2011mcmc} is highly
recommended. 

\begin{thebibliography}{99}

  \bibitem[Neal(2011)]{neal2011mcmc}
Radford M. Neal.
\newblock MCMC using Hamiltonian dynamics.
\newblock In \emph{Handbook of Markov Chain Monte Carlo}, pages 47--95.
\newblock Chapman and Hall/CRC, 2011.

%\bibitem[Girolami and Calderhead(2011)]{girolami2011riemann}
%Mark Girolami and Ben Calderhead.
%\newblock Riemann manifold Langevin and Hamiltonian Monte Carlo methods.
%\newblock \emph{Journal of the Royal Statistical Society: Series B (Statistical Methodology)}, 73(2):123--214, 2011.

%  \bibitem[Lai et~al.(2025)Lai, Song, Kim, Mitsufuji, and Ermon]{lai2025principles}
%Lai, C.-H., Song, Y., Kim, D., Mitsufuji, Y., and Ermon, S. (2025).
%\newblock The principles of diffusion models.
%\newblock \emph{arXiv preprint arXiv:2510.21890}.

%\bibitem[Glatt-Holtz et~al.(2024)]{glatt2024sacred}
%Nathan E. Glatt-Holtz, Andrew J. Holbrook, Justin A. Krometis, Cecilia F. Mondaini, and Ami Sheth.
%\newblock Sacred and Profane: from the Involutive Theory of MCMC to Helpful Hamiltonian Hacks.
%\newblock \emph{arXiv preprint arXiv:2410.17398}, 2024.

%\bibitem[Roberts and Rosenthal(2004)]{RobertsRosenthal2004}
%Roberts, G. O. and Rosenthal, J. S. (2004).
%General state space Markov chains and MCMC algorithms.
%\textit{Probability Surveys}, \textbf{1}, 20--71.

%\bibitem[Meyn and Tweedie(2009)]{MeynTweedie2009}
%Meyn, S. P. and Tweedie, R. L. (2009).
%\textit{Markov Chains and Stochastic Stability}, 2nd ed.
%Cambridge University Press.


  
%\bibitem[Chib and Greenberg(1995)]{chib1995}
%Chib, S. and Greenberg, E. (1995).
%Understanding the Metropolis--Hastings algorithm.
%\textit{The American Statistician}, \textbf{49}, 327--335.


%\bibitem[MacKay and Peto(1995)]{MacKay1995HierarchicalDirichlet}
%D.~J.~C. MacKay and L.~C. Peto,
%\newblock ``A hierarchical Dirichlet language model,''
%\newblock \emph{Natural Language Engineering}, vol.~1,
%no.~3, pp.~289--308, 1995.

%\bibitem[Rubin(1981)]{Rubin1981BayesianBootstrap}
%D.~B. Rubin,
%\newblock ``The Bayesian bootstrap,''
%\newblock \emph{The Annals of Statistics}, vol.~9,
%no.~1, pp.~130--134, 1981.
  
%\bibitem[Fisher(1934)]{Fisher1934}
%R.~A. Fisher,
%\newblock ``Probability, likelihood and quantity of information in the logic of
%uncertain inference,''
%\newblock \emph{Proceedings of the Royal Society of London. Series~A,
%Containing Papers of a Mathematical and Physical Character}, vol.~146,
%no.~856, pp.~1--8, 1934.

%\bibitem[Efron(2010)]{Efron}
%Efron, B. (2010)
%\textit{Large-scale inference: empirical Bayes methods for estimation,
%testing, and prediction}. Cambridge University Press, 2010.

%\bibitem[Shumway and Stoffer(2017)]{ShumwayStoffer}
%Shumway, R. H. and Stoffer, D. S. (2017)
%\textit{Time Series Analysis and Its Applications: With R
%  Examples}. Springer, 4th edition. 
  
%     \bibitem[Jaynes(2003)]{Jaynes} Jaynes, E. T, (2003)
%  \textit{Probability theory: the logic of science}. Cambridge
%  University Press, 2003.
 
%  \bibitem[Sinai(1992)]{Sinai} Sinai, Y. G, (1992)
%  \textit{Probability theory: an introductory course}. Springer
%  Verlag, 1992.  

%\bibitem[{deGroot(1986)}]{DeGrootBlackwell}DeGroot, M. H. (1986)
%       A conversation with David Blackwell.
%\newblock \emph{Statistical Science}, \textbf{1}, 40--53.

%         \bibitem[Mosteller(1987)]{Mosteller} Mosteller, F, (1987)
%  \textit{Fifty challenging problems in probability with
%    solutions}. Courier Corporation, 1987.

%    \bibitem[MacKay(2003)]{MacKay} MacKay, D. J. C., (2003)
%  \textit{Information theory, inference, and learning
%  algorithms}. Cambridge University Press, 2003.

%\bibitem[{Lindley(2013)}]{Lindley}Lindley, D. V. (2013)
%   \textit{Understanding Uncertainty}. John Wiley and sons, 2013.


\end{thebibliography}







\end{document}

%%% Local Variables:
%%% mode: latex
%%% TeX-master: t
%%% End:
